\section{Terms, spines and the block model}\label{sec:prelims}

\paragraph{Terms.} We write $\mathcal S$ for the closed terms built from the constant $S$
by application, associating to the left. Weak reduction $\to$ is the compatible closure
of $S\,x\,y\,z\to x\,z\,(y\,z)$, and $\twoheadrightarrow$ its reflexive--transitive
closure. Every $t\in\mathcal S$ is uniquely $S\,u_1\cdots u_m$ with $m\ge0$; the list
$\spine(t)=[u_1,\dots,u_m]$ is its \emph{spine} and the $u_i$ its \emph{items}. A term is
a \emph{head normal form} if its spine has at most two items, that is, if it is $S$,
$S\,X$ or $S\,X\,Y$; it \emph{has} a head normal form if it reduces to one. If
$\spine(t)=[x,y,z]\app\rho$, the \emph{head step} contracts the leftmost redex,
$t\to t'=x\,z\,(y\,z)\,\rho$, and
\begin{equation}\label{eq:headstep}
\spine(t')=\spine(x)\app[z,\;y\,z]\app\rho .
\end{equation}

\paragraph{The block model.} Let $B=\{0,\dots,37\}$ and let $\beta\colon B\times B\to B$
be the table of \lean{HnfS.blockTable}. Its elements are \emph{blocks}; block $0$
interprets $S$ and block $36$, written $\bot$, plays the role of ``no head normal
form''. The \emph{block} of a term is the fold $b(S)=0$, $b(u\,v)=\beta(b(u),b(v))$. The
table was found by computer search; the proofs below use only the following finite facts,
each checked by evaluation in the Lean kernel.

\begin{lemma}[table facts]\label{lem:table}
For all $a,a',a''\in B$:
\begin{enumerate}
\item[(i)] $\beta(\beta(\beta(0,a),a'),a'')=\beta(\beta(a,a''),\beta(a',a''))$
  \textup{(the $S$-law; $38^3=54{,}872$ instances)};
\item[(ii)] $\beta(0,a)\ne\bot$ and $\beta(\beta(0,a),a')\ne\bot$;
\item[(iii)] $\beta(\bot,a)=\bot$.
\end{enumerate}
\end{lemma}

\begin{lemma}[invariance]\label{lem:inv}
If $t\twoheadrightarrow t'$ then $b(t')=b(t)$.
\end{lemma}
\begin{proof}
For a single contraction of $S\,x\,y\,z$ this is Lemma~\ref{lem:table}(i) at the
blocks of $x$, $y$ and~$z$. The fold is compositional, so equality passes to contexts, and then
along the reduction.
\end{proof}

By (ii) a head normal form never has block $\bot$. A block other than $\bot$ is an
\emph{h-block}, and $t$ is \emph{h-fold} if $b(t)$ is one.

\paragraph{Walks.} The \emph{walk} of a list of terms $[u_1,\dots,u_m]$ from a block $s$
is $q_0=s$, $q_i=\beta(q_{i-1},b(u_i))$. From $s=0$ the walk of $\spine(t)$ ends at
$b(t)$, since $t=S\,u_1\cdots u_m$. By Lemma~\ref{lem:table}(iii), a walk that ends at an
h-block visits only h-blocks.
